\documentclass[11pt]{article}
\usepackage[a4paper, margin=1in]{geometry}
\usepackage{amsmath}
\usepackage{amsfonts}
\usepackage{amssymb}
\usepackage{enumitem}

\title{Fórmulas de Análisis de Créditos}
\author{}
\date{}

\begin{document}
\maketitle

\section{Cuota}

\begin{center}
    \textbf{Cuota} = \textbf{Interés} + \textbf{Amortización}
\end{center}

La fórmula para calcular la renta (cuota) de un préstamo es la siguiente:

\begin{equation}
    R = P \bigg[ \dfrac{(1 + i)^{n} \cdot i}{(1 + i)^{n} - 1}\bigg]
\end{equation}

Donde:
\begin{itemize}[label={--}]
    \item \textbf{P}: Valor o cantidad de dinero en el momento denominado "presente" o "tiempo cero".
    \item \textbf{R}: Renta o cuota a pagar.
    \item \textbf{i}: Tasa de interés efectiva por periodo de interés (interés compuesto).
    \item \textbf{n}: Número de periodos de capitalización.
\end{itemize}

\section{Interés}

\begin{center}
    \textbf{Intereses} += \textbf{Saldo} $\cdot$ \textbf{Interés}
\end{center}

\begin{equation}
    TEM = (1 + TEA)^{\frac{1}{12}} - 1
\end{equation}

Donde:
\begin{itemize}[label={--}]
    \item  \textbf{TEM}: Tasa de efectiva mensual.
    \item  \textbf{TEA}: Tasa de efectiva anual.
    \item \textbf{I}: Interés generado.
    \item \textbf{S}: Saldo o monto del capital en el periodo.
    \item \textbf{i}: Tasa de interés por periodo (en decimal).
\end{itemize}

\section{Amortización}

\begin{center}
    \textbf{Amortización} = \textbf{Cuota} - \textbf{Intereses}
\end{center}


\end{document}
